\chapter{CONCLUSIONES Y RECOMENDACIONES}
\label{cap:conclusiones}

\section{Conclusiones}
\label{sec:conclusiones}

El presente proyecto de grado cumplió el objetivo general de desarrollar un
sistema inteligente para el entrenamiento en gimnasios que proporciona una guía
estructurada y un seguimiento personalizado de los ejercicios. El producto es
una aplicación Android construida con React Native, Expo y TypeScript sobre
los servicios de Firebase, organizada en una arquitectura cliente--servidor con
capas Modelo--Vista--Controlador y desarrollada con la metodología Mobile-D.

Respecto a cada objetivo específico, se concluye lo siguiente:

\begin{enumerate}[label=\textbf{OE\arabic*.}, leftmargin=3.2em]
    \item \textbf{Guiar sobre el uso correcto de máquinas.} Se construyó un
    catálogo de 35 tipos de máquinas, con 288 modelos comerciales como
    variantes, en el que cada ficha incluye descripción, músculos trabajados,
    guía de uso paso a paso, errores comunes y pautas de prevención de
    lesiones redactadas como contenido propio. El catálogo se consulta con
    búsqueda tolerante a tildes y filtros, y se complementa con un escáner que
    superpone la ficha de la máquina sobre la imagen en vivo de la cámara.
    \item \textbf{Instruir en la ejecución correcta de ejercicios.} Cada
    máquina presenta al menos un ejercicio con su animación de ejecución, su
    músculo objetivo y su nivel de dificultad, con un total de 95 animaciones
    obtenidas por la vía oficial de su proveedor.
    \item \textbf{Diseñar rutinas personalizadas.} Se implementó un generador
    de rutinas basado en reglas cuya base de conocimiento traduce las pautas
    de prescripción del ejercicio en reglas de frecuencia, de variables de
    carga y de división semanal. Las pruebas sobre las doce combinaciones de
    objetivo y nivel confirmaron que cada rutina generada respeta esas reglas.
    La elección de un enfoque basado en conocimiento, en lugar de uno de
    aprendizaje automático, permitió obtener recomendaciones explicables sin
    disponer de un conjunto de datos de entrenamiento.
    \item \textbf{Construir herramientas de planificación.} El planificador
    de rutinas personalizadas permite organizar la semana día por día, ajustar
    series, repeticiones y descansos, mantener varias rutinas con una activa y
    agregar ejercicios directamente desde la ficha de una máquina.
    \item \textbf{Registrar la ingesta de alimentos.} Dentro del alcance
    definido, que limita este objetivo al control calórico mediante una
    calculadora de IMC, se implementó el cálculo y la clasificación del IMC
    según la OMS con el registro histórico de cada medición. El modelo de
    datos y las reglas de seguridad ya contemplan el registro de ingesta, lo
    que facilita su ampliación.
    \item \textbf{Implementar el seguimiento del progreso físico.} El usuario
    registra su peso y su estatura, el sistema calcula el IMC y el historial
    se sincroniza en tiempo real en todos sus dispositivos.
    \item \textbf{Generar reportes de progreso.} El reporte presenta la
    evolución del peso en un gráfico, la variación acumulada y la categoría de
    IMC actual, calculados al vuelo a partir del historial para reflejar
    siempre los datos vigentes.
\end{enumerate}

En el plano técnico, se concluye además que:

\begin{itemize}
    \item La separación en capas resultó determinante para la mantenibilidad:
    al aislar el acceso a Firebase en los repositorios y servicios, cambios
    como el paso de colecciones raíz a subcolecciones por usuario, o la
    sustitución de la fuente del catálogo, se resolvieron sin modificar las
    pantallas.
    \item Las reglas de seguridad del servidor, que validan propiedad, tipos y
    rangos, son indispensables en una arquitectura sin servidor propio, porque
    las credenciales del cliente son públicas por diseño.
    \item El diseño del reconocimiento automático como un punto de extensión
    con un contrato definido permite incorporar en el futuro un clasificador
    de imágenes sin afectar al resto del sistema.
\end{itemize}

\section{Recomendaciones}
\label{sec:recomendaciones}

A partir de la experiencia del desarrollo se recomienda, como trabajo futuro:

\begin{itemize}
    \item \textbf{Reconocimiento automático de máquinas con TensorFlow Lite.}
    Construir un conjunto de fotografías propias de las 35 máquinas, tomadas
    en gimnasios con autorización, entrenar un clasificador por transferencia
    de aprendizaje sobre una arquitectura ligera como MobileNet e implementarlo
    en la función \archivo{reconocerMaquina}. No deben usarse las animaciones
    de WorkoutX para el entrenamiento, porque sus condiciones de uso lo
    prohíben sin permiso escrito.
    \item \textbf{Realidad aumentada con anclaje tridimensional.} Evaluar
    ARCore para anclar indicaciones sobre la máquina real (puntos de ajuste,
    trayectoria del movimiento), como evolución de la superposición basada en
    cámara.
    \item \textbf{Progresión de la carga.} Ampliar la base de conocimiento del
    generador con reglas de sobrecarga progresiva que ajusten la rutina según
    el historial del usuario.
    \item \textbf{Diario de alimentación.} Incorporar la interfaz de registro
    de ingesta sobre el modelo ya preparado y mostrar el historial de IMC en
    la pantalla de nutrición.
    \item \textbf{Pruebas automatizadas.} Instalar \archivo{jest-expo} y
    React Native Testing Library para convertir las pruebas de la lógica de
    negocio en una suite ejecutable en integración continua, y probar las
    reglas con el emulador de Firebase.
    \item \textbf{Validación con usuarios.} Realizar un estudio de usabilidad
    con usuarios reales de gimnasio, por ejemplo con el cuestionario SUS, y
    medir el efecto de la aplicación en la permanencia de los usuarios nuevos.
    \item \textbf{Publicación.} Antes de publicar la aplicación en Google
    Play, obtener la autorización del proveedor de las fotografías del
    catálogo y verificar las condiciones de uso de las animaciones.
\end{itemize}
